% C4 Architecture Visualization Documentation Suite | Original educational content
% Standalone LaTeX root: no external style file, no network, no shell escape.
\documentclass[10pt,letterpaper]{article}
\usepackage[T1]{fontenc}
\usepackage[utf8]{inputenc}
\usepackage{lmodern,microtype}
\usepackage[letterpaper,margin=0.82in,headheight=15pt,headsep=18pt,footskip=26pt]{geometry}
\usepackage[table]{xcolor}
\usepackage{booktabs,tabularx,array,longtable}
\usepackage{enumitem}
\usepackage{fancyhdr,lastpage}
\usepackage{titlesec}
\usepackage{hyperref,xurl}
\usepackage{needspace}
\usepackage{listings}
\usepackage{amssymb}
\definecolor{C4Navy}{HTML}{17364B}
\definecolor{C4Teal}{HTML}{137C85}
\definecolor{C4Gray}{HTML}{526775}
\definecolor{C4Light}{HTML}{EEF4F5}
\definecolor{C4Line}{HTML}{CAD8DE}
\color{C4Navy}
\hypersetup{colorlinks=true,linkcolor=C4Teal,urlcolor=C4Teal,citecolor=C4Teal}
\setlength{\parindent}{0pt}
\setlength{\parskip}{5pt}
\setlength{\emergencystretch}{2em}
\raggedbottom
\setlist[itemize]{leftmargin=1.4em,topsep=3pt,itemsep=3pt}
\setlist[enumerate]{leftmargin=1.6em,topsep=3pt,itemsep=3pt}
\titleformat{\section}{\large\bfseries\sffamily\color{C4Navy}}{\thesection}{0.6em}{}
\titleformat{\subsection}{\normalsize\bfseries\sffamily\color{C4Teal}}{\thesubsection}{0.6em}{}
\titlespacing*{\section}{0pt}{16pt}{6pt}
\titlespacing*{\subsection}{0pt}{9pt}{4pt}
\pagestyle{fancy}
\fancyhf{}
\fancyhead[L]{\scriptsize\sffamily C4 / ARCHITECTURE VISUALIZATION}
\fancyhead[R]{\scriptsize\sffamily REFERENCE / EDUCATION}
\fancyfoot[L]{\scriptsize\color{C4Gray}ORIGINAL ENGINEERING SUITE / 2026}
\fancyfoot[R]{\scriptsize\thepage\ / \pageref*{LastPage}}
\renewcommand{\headrulewidth}{0.4pt}
\newcolumntype{Y}{>{\raggedright\arraybackslash}X}
\newcommand{\RuleLine}{\par\noindent\color{C4Line}\rule{\linewidth}{0.45pt}\color{C4Navy}\par}
\newcommand{\WritingSpace}[1]{\par\vspace{2pt}\noindent\fcolorbox{C4Line}{white}{\parbox[t][#1][t]{0.94\linewidth}{\hspace{1pt}}}\par\vspace{7pt}}
\newcommand{\Field}[1]{\par\textbf{#1}\par\WritingSpace{0.57in}}
\newcommand{\Prompter}[1]{\par{\small\color{C4Teal}\textbf{Authoring guidance:} #1}\par}
\newcommand{\SourceNote}[1]{\par\smallskip\noindent\colorbox{C4Light}{\parbox{0.94\linewidth}{\footnotesize\textbf{PROVENANCE.} #1}}\par\medskip}
\lstset{basicstyle=\ttfamily\footnotesize,breaklines=true,columns=fullflexible,showstringspaces=false,frame=single,rulecolor=\color{C4Line}}
\newcommand{\DocStart}[3]{%
\begin{document}\begin{titlepage}
\vspace*{0.55in}
{\sffamily\bfseries\small\color{C4Teal}#1}\par\vspace{0.24in}
{\sffamily\bfseries\fontsize{26}{31}\selectfont\color{C4Navy}#2}\par\vspace{0.18in}
{\large\color{C4Gray}#3}\par\vspace{0.4in}
\textcolor{C4Teal}{\rule{\linewidth}{1.15pt}}\par\vspace{0.24in}
\textbf{Document control}\par
Edition 1.0\quad|\quad 30 October 2026\quad|\quad Original educational engineering material\par
\vspace{0.17in}
\textbf{Source boundary}\par
Simon Brown, \emph{The C4 Model: Visualizing Software Architecture} (O'Reilly, 2026), is the primary study source. The suite's exercises, forms, crosswalks, process rules, sample systems, and scoring models are independently authored supplements; they are not represented as the author's prescriptions.\par
\vfill
{\small\color{C4Gray}Not an ISO/IEC/IEEE 42010 conformance certificate. Not a substitute for the published standard or for review of deployed systems.}\par
\end{titlepage}
}
\DocStart{COMPREHENSIVE TECHNICAL REFERENCE}{C4 Architecture Visualization Engineering Manual}{Semantics, executable diagram standards, traceability, crosswalks and governance}
\tableofcontents\clearpage
\SourceNote{Part I presents original instructional synthesis of Simon Brown (2026). Parts II--IV are original enterprise engineering extensions that integrate C4 with the repository, PlantUML, ISO/IEC/IEEE 42010, UML, and Views and Beyond. The extensions are not claimed to be book prescriptions or normative standard requirements.}
\section{Executive engineering profile}
The C4 model should be treated as an abstraction and communication vocabulary rather than a mandated notation, diagram count, security proof or deployment inventory. Every architecture illustration should state (1) the system/area under discussion, (2) the viewpoint and concern, (3) the abstraction level, (4) which assertions are observed versus proposed, and (5) who reviewed the supporting evidence.

\begin{tabularx}{\linewidth}{@{}YYY@{}}\toprule
\textbf{Responsibility} & \textbf{Appropriate artifact} & \textbf{What it does not prove} \\ \midrule
Describe context & C4 system context diagram & Internal decomposition or security correctness \\
Describe software runtime building blocks & C4 container diagram & Deployment hosts, node topology or every code package \\
Describe internal responsibilities & C4 component diagram & Complete source structure \\
Explain runtime behavior & Dynamic diagram, optionally UML sequence & All possible paths or guaranteed execution ordering \\
Describe placement & Environment-specific deployment diagram & Running state without runtime evidence \\
Record consequential choice & Architecture decision record & The complete architecture description \\
Organize stakeholder coverage & Architecture description with viewpoints and views & Automatic conformance with a standard \\
\bottomrule\end{tabularx}

\section{Part I — Book-aligned C4 abstractions}
\subsection{Hierarchy and invariants}
At the system context zoom, a software system appears as one unit. A container view opens that boundary and shows major applications and data stores. At component zoom, one container is selected and internally decomposed. The optional code view shows implementation details that matter to a particular question. Supporting diagrams answer runtime, deployment, and cross-system landscape questions rather than creating a fifth zoom. A box is therefore never interpreted from shape alone.

\begin{itemize}
\item Containment: a component belongs to a container, and a container belongs to a software system in a defined model.
\item Identity: the same logical system or container must retain the same name/ID between views.
\item Relation: direction and action text must make sense to a reader without the author narrating.
\item Perspective: each view has one explicit abstraction and should not silently combine levels.
\item Provenance: technology labels and deployment claims are evidence-based or clearly marked as proposed/unknown.
\end{itemize}
\subsection{Diagram selection decision procedure}
\begin{tabularx}{\linewidth}{@{}YYY@{}}\toprule
\textbf{Reader question} & \textbf{Select} & \textbf{Intentionally exclude} \\ \midrule
Who uses the system and what integrates with it? & Context & Internal apps and database details \\
What runs and stores data inside the system? & Container & Hosts and node placements \\
How is a complex application partitioned? & Component & Exhaustive class inventory \\
How does this algorithm or API type work? & Code & Unrelated classes \\
How does one business scenario execute? & Dynamic & All other scenarios \\
Where is the application deployed for this environment? & Deployment & Assumed equivalence of all environments \\
How do multiple systems relate across a business domain? & Landscape & Single-system internals \\
\bottomrule\end{tabularx}

\subsection{Book and tool separation}
Brown discusses four zoom levels and three supporting diagram types. PlantUML is an optional diagram-as-code implementation, not a required C4 notation. The repository house style uses C4-PlantUML to maintain consistency, while typed macros and syntax are tool-specific. References to a given library macro are tooling guidance, not assertions about the definition of software architecture.

\section{Part II — Diagram specification and authoring contract}
\subsection{System context — content and acceptance}
\begin{tabularx}{\linewidth}{@{}YY@{}}\toprule
\textbf{Facet} & \textbf{Requirement} \\ \midrule
Primary question & What is the system of interest and its environment? \\
Minimum content & A single software system, roles and external software systems, direction and interaction purpose. \\
Authoring boundary & Use it as a boundary agreement: include direct dependencies but no internal application, database or deployment node. \\
Independent interpretation test & The business sponsor can identify the system, distinguish external systems and explain its business purpose. \\
\bottomrule\end{tabularx}

\subsection{Container — content and acceptance}
\begin{tabularx}{\linewidth}{@{}YY@{}}\toprule
\textbf{Facet} & \textbf{Requirement} \\ \midrule
Primary question & What independently running applications and data stores make up the system? \\
Minimum content & An enclosing software-system boundary, application/data store types, significant technology choices and interactions. \\
Authoring boundary & Map externally visible context interactions to internal origins/receivers; label significant protocols and state unknown facts. \\
Independent interpretation test & Platform lead can identify ownership, runtime building blocks, dependencies and data stores. \\
\bottomrule\end{tabularx}

\subsection{Component — content and acceptance}
\begin{tabularx}{\linewidth}{@{}YY@{}}\toprule
\textbf{Facet} & \textbf{Requirement} \\ \midrule
Primary question & How is a particular container internally structured? \\
Minimum content & Selected container, architecturally significant components, responsibilities, interfaces and clients. \\
Authoring boundary & Give source mapping for nontrivial components; do not imply a one-to-one correspondence between teams and components. \\
Independent interpretation test & Maintainer can locate the responsible implementation and understand critical dependencies. \\
\bottomrule\end{tabularx}

\subsection{Code — content and acceptance}
\begin{tabularx}{\linewidth}{@{}YY@{}}\toprule
\textbf{Facet} & \textbf{Requirement} \\ \midrule
Primary question & Which implementation relationships need explanation? \\
Minimum content & Classes/interfaces/functions important to a selected question, at a named code revision. \\
Authoring boundary & Use code-level semantics appropriate to the claimed notation; retire mechanically generated clutter. \\
Independent interpretation test & Developer can connect diagram relations to actual source constructs. \\
\bottomrule\end{tabularx}

\subsection{Dynamic — content and acceptance}
\begin{tabularx}{\linewidth}{@{}YY@{}}\toprule
\textbf{Facet} & \textbf{Requirement} \\ \midrule
Primary question & How does one scenario proceed? \\
Minimum content & Consistent participants, numbered/ordered interactions, request/response, failure/retry where needed. \\
Authoring boundary & Use the declared C4 zoom and static aliases; distinguish ordinary success from exceptional behavior. \\
Independent interpretation test & Reviewer can follow one scenario and inspect every interaction as evidence. \\
\bottomrule\end{tabularx}

\subsection{Deployment — content and acceptance}
\begin{tabularx}{\linewidth}{@{}YY@{}}\toprule
\textbf{Facet} & \textbf{Requirement} \\ \midrule
Primary question & Where do known container instances execute? \\
Minimum content & Named environment, nodes, network zones, container instances, external service connections. \\
Authoring boundary & Map logical container aliases to each instance; do not label hoped-for redundancy as verified reliability. \\
Independent interpretation test & SRE can locate instances, ingress, data services and trust boundaries. \\
\bottomrule\end{tabularx}

\subsection{Landscape — content and acceptance}
\begin{tabularx}{\linewidth}{@{}YY@{}}\toprule
\textbf{Facet} & \textbf{Requirement} \\ \midrule
Primary question & How do related software systems fit inside an estate or capability? \\
Minimum content & Enterprise/domain scope, named systems, users, integration relations and owners. \\
Authoring boundary & Do not duplicate internal container detail; use a connected system catalog. \\
Independent interpretation test & Portfolio stakeholders can identify system ownership and integration seams. \\
\bottomrule\end{tabularx}

\section{View packet contract (supplementary)}
A production-ready diagram packet is more than a rendered image. The packet should let a reviewer identify its purpose, check its semantics, follow key dependencies, and reproduce the render.

\begin{itemize}
\item View identifier, model version and system-of-interest identifier.
\item Viewpoint identifier, stakeholders and named concerns.
\item Source file with a pinned rendering toolchain reference and exported image, where appropriate.
\item Element register: IDs, types, descriptions, owner or authoritative source.
\item Relationship register: initiating ID, receiving ID, verb, protocol and data classification where relevant.
\item Correspondence links to adjacent zoom levels, behavioral scenarios or infrastructure evidence.
\item Review log with authorship, reviewer, discovered ambiguity, decision and revalidation requirement.
\end{itemize}
\section{Original fictional case study: Civic Permit Portal}
This is a teaching model, not a verified production architecture. A resident submits a building-permit application; city reviewers perform manual review; an external identity provider authenticates the caller; a payment provider authorizes fees; a GIS integration supplies parcel data.

\begin{tabularx}{\linewidth}{@{}YYY@{}}\toprule
\textbf{Element ID} & \textbf{C4 abstraction} & \textbf{Responsibility / evidence status} \\ \midrule
SYS-PERMIT & Software system & Permit submission, status and review; illustrative \\
PER-RESIDENT & Person & Applicant and status viewer \\
SYS-IDP & External system & Authentication integration; conceptual \\
SYS-FEES & External system & Fee authorization; conceptual \\
CNT-WEB & Container & React UI; selected teaching technology \\
CNT-API & Container & FastAPI domain API; selected teaching technology \\
CNT-WORKER & Container & Background review processor; selected teaching technology \\
CNT-STORE & Container (data store) & PostgreSQL permit record database; teaching choice \\
CMP-FLOW & Component & Permit workflow policy implementation; illustrative \\
\bottomrule\end{tabularx}

\subsection{Trace a single cross-view relationship}
At context zoom, SYS-PERMIT requests fee authorization from SYS-FEES. At container zoom, CNT-API requests payment authorization from SYS-FEES. At dynamic zoom, the payment step occurs between form acceptance and record persistence. At deployment zoom, the API instance originates egress traffic from a production execution node. At implementation zoom, CMP-PAYMENT maps to a package and integration contract. This chain is illustrative; actual checks must use a deployed service, real code or an explicitly marked target state.

\section{Part III — Traceability, architecture descriptions and standards integration}
\subsection{ISO/IEC/IEEE 42010:2022 architecture-description lens}
The published standard distinguishes an architecture from an architecture description and gives requirements for describing entities of interest, viewpoints, model kinds and relationships among description concepts. Its official scope does not mandate a particular drawing notation, authoring tool or C4 diagram taxonomy. These suite fields form a proposed documentation profile. Teams must inspect the actual normative requirements before claiming conformance.

\begin{tabularx}{\linewidth}{@{}YYY@{}}\toprule
\textbf{Concern-management concept} & \textbf{Proposed C4 implementation} & \textbf{Evidence to retain} \\ \midrule
Stakeholder and concern & Concern record linked to a selected C4 viewpoint & Concern origin and acceptance question \\
Architecture viewpoint & Reusable view specification from viewpoints/ & Intended scope, model kinds and interpretation rules \\
Architecture view & Versioned C4 diagram record & View ID and model revision \\
Model kind & Context, container, component, code, dynamic, deployment or landscape notation definition & Visible abstraction rules and diagram key \\
Correspondence & Cross-view workbook entry & From/to IDs, rationale and validation status \\
Architecture decision & Separate ADR in existing decisions collection & Context, alternatives, rationale and consequences \\
\bottomrule\end{tabularx}

\subsection{Views and Beyond complement}
Views and Beyond organizes documentation concerns around module, component-and-connector and allocation viewtypes. A C4 container diagram can reveal runtime building blocks yet often cannot express a complete source-code module-use hierarchy. A C4 deployment can illustrate placement yet may not satisfy all allocation questions. The architecture package should include complementary views where concerns demand them; these are not one-to-one translations.

\begin{tabularx}{\linewidth}{@{}YYY@{}}\toprule
\textbf{Views and Beyond emphasis} & \textbf{Potential companion artifact} & \textbf{Do not equate} \\ \midrule
Module decomposition, layering, uses & Module dependencies, source package map, code ownership & A team structure with a C4 component hierarchy \\
Runtime component and connector structure & C4 containers, component relations, protocol contract & Every C4 static relationship with observed runtime behavior \\
Allocation to infrastructure, code base or teams & C4 deployment plus host-to-instance and repo mappings & C4 deployment with all allocation concerns \\
Interfaces and behavior & UML sequence, interface contract, state model & A readable drawing with complete interface semantics \\
\bottomrule\end{tabularx}

\subsection{UML relationship}
UML 2.5.1 provides detailed structural and behavioral constructs with their own semantics. A UML sequence diagram can express ordering, lifelines and alternatives beyond a simplified C4 dynamic view. A UML class diagram can capture a narrow set of code-level structural relationships. A mixed document should clearly identify where a figure is C4-style versus normatively interpreted UML. A correct PlantUML render does not verify the semantics or completeness of either.

\subsection{Traceability schema and validation steps}
\begin{lstlisting}
# Example evidence records: suite design, not C4 syntax
concern_id: CONC-004
stakeholder_id: STAKE-REVIEWER
viewpoint_id: VP-C4-DYNAMIC
view_id: VIEW-PERMIT-SUBMIT
model_revision: example-commit-sha
related_elements: [CNT-WEB, CNT-API, SYS-FEES]
evidence:
  - kind: code
    location: src/services/permit/workflow.py
    observation: proposed_example
  - kind: runtime
    location: none
    observation: not_verified
review_state: draft
\end{lstlisting}
\begin{itemize}
\item Validate that every referenced concern, viewpoint, view and element exists in a canonical registry.
\item Check containment edges: component-to-container, container-to-system and instance-to-container.
\item Ensure edge endpoints exist, have meaningful verbs and use IDs consistent with neighboring views.
\item Ensure actual/observed assertions point to evidence revisions and observed dates.
\item Generate a discrepancy report; human reviewers adjudicate ownership, meaning, risk and acceptability.
\end{itemize}
\section{Part IV — Operate the documentation as an engineering product}
\subsection{Canonical PlantUML source pattern}
\begin{lstlisting}
@startuml civic-permit-context
!include ../../../../../tooling/styles/plantuml/c4/context-diagram-style.iuml
C4_TITLE(Civic Permit Portal, system context)
Person(applicant, "Applicant", "Files applications")
System(portal, "Permit Portal", "Handles permits")
System_Ext(identity, "Identity Provider", "Authenticates users")
Rel(applicant, portal, "Submits an application", "HTTPS")
Rel(portal, identity, "Authenticates callers", "OIDC")
SHOW_LEGEND()
@enduml
\end{lstlisting}
In this repository, the single house-style include brings pinned C4-PlantUML macros and shared visual tokens. Do not add local palette overrides or retired skinparam directives. The relative path assumes the file is at the delivered examples/ directory depth; relocate it only after adjusting the include.

\subsection{Pipeline and semantic gates}
\begin{tabularx}{\linewidth}{@{}YYY@{}}\toprule
\textbf{Stage} & \textbf{Machine-checkable portion} & \textbf{Requires human judgment} \\ \midrule
Source intake & File naming, expected root and paired-form presence & Whether this artifact should exist \\
PlantUML lint & House-style include, forbidden local styling, syntax of known patterns & Whether the diagram expresses the architecture \\
Render & Diagram source compiles under pinned renderer & Visual legibility and interpretation \\
Model crosscheck & Aliases resolve, relationships have endpoints, references link & Whether dependency is valid and secure \\
Evidence drift & Changed IaC/code paths or runtime inventory references & Whether an architecture update is material \\
Approval & Reviewer IDs, completed fields and thresholds & Risk acceptance and standards conformity \\
\bottomrule\end{tabularx}

\subsection{Security and operational risk lens}
\begin{itemize}
\item Mark crossings from untrusted callers to public APIs and from application workloads to third-party services.
\item Inspect identity assertions separately from access-control enforcement; a line called authenticates is not proof of authorization.
\item Track sensitive records through data stores, queues and exports rather than assuming encryption from technology labels.
\item For message-driven flows, document delivery semantics, deduplication, retries, dead-letter handling and consumer ownership where important.
\item For deployment, prove public/private placement, workload identities and network permissions from deployed cloud inventory, not merely a Terraform intent file.
\end{itemize}
\subsection{Review score and hard-stop policy}
The attached scorecards allocate weighted criteria, each rated 0 through 4. Compute a weighted score as the sum over criteria of weight multiplied by rating divided by 4. This is a proposed scoring method for consistent review, not a book or ISO standard grading scale. A critical blocker, such as a known unsupported external trust relationship, prevents approval regardless of total score.

\begin{lstlisting}
score_percent = sum(weight_percent_i * rating_i / 4)
release_eligible = (score_percent >= local_threshold) and no_critical_blockers
\end{lstlisting}
\subsection{Operating cadence and ownership}
\begin{tabularx}{\linewidth}{@{}YYY@{}}\toprule
\textbf{Trigger} & \textbf{Expected review} & \textbf{Artifact owner} \\ \midrule
New external integration & Context and container edges; associated security view & System architect and AppSec \\
Service decomposition changes & Container and component IDs; affected ADRs & Engineering lead \\
New cloud topology & Deployment nodes/instances and runtime evidence & Platform/SRE \\
New event choreography & Dynamic scenarios and failure modes & Application engineer/SRE \\
Quarterly review & Age of evidence and changed dependencies & Architecture steward \\
\bottomrule\end{tabularx}

\section{Implementation acceptance checklist}
\begin{itemize}
\item The suite is installed at the domain-appropriate repository path; README and catalogs link to it.
\item The chosen diagram set is justified by stakeholder questions, not diagram completeness pressure.
\item Every published view has a versioned source and a legend that resolves notational ambiguity.
\item Cross-view mappings have no unexplained missing endpoints.
\item Every claim of production placement or security posture is backed by actual evidence or marked unverified.
\item Reusable authoring Template documents have separate print-friendly Blank files.
\item Diagram lint, LaTeX build, independent diagram review and disposition all pass.
\item Any compliance claim is assessed separately against the applicable published normative standard.
\end{itemize}
\section{Source register and further study}
\begin{itemize}
\item Primary source: Brown, The C4 Model: Visualizing Software Architecture (O’Reilly, 2026), Chapters 1–12 and Appendix, user-supplied file.
\item Official C4 model, glossary and review checklist: \url{https://c4model.com/diagrams} and \url{https://c4model.com/diagrams/checklist}.
\item ISO/IEC/IEEE 42010:2022, official ISO catalogue: \url{https://www.iso.org/standard/74393.html}.
\item OMG UML 2.5.1, official specification: \url{https://www.omg.org/spec/UML/2.5.1}.
\item Clements et al., Documenting Software Architectures: Views and Beyond, Second Edition, Addison-Wesley.
\item C4-PlantUML reference: \url{https://plantuml-stdlib.github.io/C4-PlantUML/}.
\item Repository style contract: tooling/plantuml/README.md and tooling/styles/plantuml/c4/.
\end{itemize}
\end{document}
